% chapters/chapter3.tex

\chapter{Methodology}\label{chap:methodology}

\section{Chapter Overview and Implemented Scope}
This chapter describes the methodology that was actually implemented, trained, and evaluated in the current prototype. The core experimental object is the \textbf{assistant-output watermarking branch}: a waveform-domain speech watermark that should remain detectable under common post-processing, keep the false-positive rate low, and recover its payload from short audio prefixes. The complementary \textbf{user-side deepfake/replay detector} remains outside the present implementation and is therefore treated as a systems-level extension rather than as an experimentally validated branch.

The practical objective is narrower than the broader middleware vision stated in Chapter~1. In the implemented scope, the thesis studies a \textbf{localized, prefix-aware watermark for 16~kHz speech on LibriSpeech}. This scope is still meaningful for a conversational setting, because a deployed verifier does not want to wait for the full utterance whenever an earlier and sufficiently reliable decision is already possible. Accordingly, the methodology below is organized around four concrete questions:
\begin{enumerate}[leftmargin=1.25em]
  \item can a small payload be embedded and decoded reliably at all?
  \item can detection remain localized instead of collapsing into a trivial always-positive decision?
  \item can payload verification happen early enough for streaming use?
  \item how does the resulting checkpoint behave under a realistic, if still modest, robustness benchmark?
\end{enumerate}

To answer these questions, the implementation uses a staged curriculum. Rather than optimizing the final localized low-latency system from scratch, the training pipeline first learns the bit channel, then calibrates positive versus negative detection, then introduces localization, and finally fine-tunes for prefix decoding. In practice this curriculum was not a stylistic preference but a stability requirement: direct end-to-end optimization of localization, payload recovery, and short-prefix decoding was unreliable in early runs.

\section{Experimental Data and Split Strategy}
\subsection{Dataset actually used in the implemented system}
Although earlier drafts discussed a broader multi-domain benchmark, the current implementation is trained and evaluated on \textbf{LibriSpeech} only. This decision reduced experimental variance and made stage-wise iteration reproducible during model development. LibriSpeech provides clean read English speech at 16~kHz together with a widely used directory structure and established train/development/test partitions~[58].

In the current code base, manifests are created directly from the extracted LibriSpeech tree. The main manifests are:
\begin{itemize}[leftmargin=1.25em]
  \item \texttt{manifests/librispeech\_train.csv},
  \item \texttt{manifests/librispeech\_val.csv},
  \item \texttt{manifests/librispeech\_test.csv}.
\end{itemize}
For faster curriculum runs, reduced manifests are also used:
\begin{itemize}[leftmargin=1.25em]
  \item \texttt{manifests/train\_2k.csv} for a 2,000-file training subset,
  \item \texttt{manifests/val\_200.csv} for a 200-file validation subset.
\end{itemize}
The reduced manifests accelerate stage transitions, while the full LibriSpeech test manifest is retained for final robustness evaluation.

\subsection{Preprocessing as implemented}
The preprocessing pipeline is intentionally simple and reflects the behavior of the code rather than an idealized design. For each file, the loader:
\begin{itemize}[leftmargin=1.25em]
  \item converts audio to mono,
  \item clips amplitudes to $[-1,1]$ if needed,
  \item resamples the signal to 16~kHz,
  \item forms a fixed-length segment by cropping or zero-padding.
\end{itemize}
No separate silence filtering or RMS normalization is applied in the reported experiments. During training, long files are sampled with random crops. During validation and testing, cropping is deterministic, which reduces evaluation noise and makes stage-to-stage comparisons easier to interpret.

Two segment lengths are used in the curriculum:
\begin{itemize}[leftmargin=1.25em]
  \item \textbf{2~s segments} in the bridge and detector-calibration stages,
  \item \textbf{4~s segments} in the localized and latency-aware stages.
\end{itemize}
This progression mirrors the actual development path: the shorter segments made it easier to stabilize the bit channel early, while the longer segments later provided room for localization holes and prefix-aware evaluation.

\section{Watermarking Formulation}
Let $x \in \mathbb{R}^{N}$ denote a clean waveform segment with $N$ samples. The encoder receives $x$ and a binary payload $m \in \{0,1\}^{B}$ and produces a watermarked waveform
\begin{equation}
  y = \mathrm{clip}(x + \delta, -1, 1),
\end{equation}
where $\delta$ is a learned additive residual and \texttt{max\_delta} bounds the residual magnitude. In the strongest stable configuration reported here, the payload size is \textbf{4 bits}. Earlier attempts with larger payloads were less stable and therefore are not presented as the main result.

The detector processes an observed waveform $\tilde{y}$ and predicts:
\begin{enumerate}[leftmargin=1.25em]
  \item frame-level watermark-presence logits $\hat{p}_{1:F}$,
  \item payload logits $\hat{m} \in \mathbb{R}^{B}$.
\end{enumerate}
Unlike an earlier draft design that decoded bits frame by frame, the implemented detector performs \emph{masked statistics pooling}: it extracts frame-level features, pools them over the predicted or known watermark region, and only then predicts the payload. This was a decisive architectural simplification for stable training.

\subsection{Detection rule}
The current inference rule uses a threshold on the fraction of frames whose presence probability exceeds a scalar threshold $\tau$:
\begin{equation}
  \widehat{\mathrm{present}} = \mathbb{1}\!\left[\frac{1}{F}\sum_{f=1}^{F}\mathbb{1}\big(\sigma(\hat{p}_f) \ge \tau\big) > \rho\right],
\end{equation}
where $\sigma(\cdot)$ is the sigmoid function and $\rho$ is the minimum positive-frame fraction. In the reported configs, $\tau = 0.5$ and $\rho = 0.05$.

\subsection{Localization mask}
When the watermark is not present over the full segment, the system maintains a binary mask $s \in \{0,1\}^{F}$ over detector frames. In the localized stages, watermark coverage is created by embedding the watermark almost everywhere and then carving one or more short \emph{holes}. This design is operationally useful for streaming verification: it trains the presence head to reason about where the watermark exists, not merely whether some watermark evidence appears somewhere in the segment.

\section{Implemented Model Architecture}
\subsection{Encoder}
The encoder is a \textbf{message-conditioned dilated residual network} operating directly on the waveform. The implementation uses:
\begin{itemize}[leftmargin=1.25em]
  \item an input projection convolution,
  \item seven conditioned residual blocks with dilations $[1,2,4,8,16,32,64]$,
  \item group normalization and GELU activations,
  \item a final 1D convolution followed by \texttt{tanh} scaling to keep the residual bounded.
\end{itemize}
Each residual block is conditioned on the payload through an affine transformation of hidden activations. This lets the message influence multiple temporal scales instead of only the input layer. With the current hyperparameters, the receptive field is 769 samples, or about 48~ms at 16~kHz.

\begin{lstlisting}[language=Python,caption={Simplified encoder forward pass from the implemented model.},label={lst:encoder_forward}]
def forward(self, x, bits):
    msg = self.msg_fc(bits)
    h = self.in_proj(x)
    for block in self.blocks:
        h = block(h, msg)
    h = F.gelu(self.out_norm(h))
    residual = torch.tanh(self.out_conv(h)) * self.max_delta
    y = torch.clamp(x + residual, -1.0, 1.0)
    return y, residual
\end{lstlisting}

\subsection{Detector}
The detector is also waveform-based. It consists of:
\begin{itemize}[leftmargin=1.25em]
  \item a downsampling 1D CNN backbone,
  \item a \textbf{presence head} that predicts frame-level watermark logits,
  \item a \textbf{bit context module} that refines backbone features,
  \item a \textbf{masked statistics pooling} block,
  \item a final MLP that predicts the payload bits.
\end{itemize}
The detector downsamples by a factor of 16, which preserves enough temporal resolution for localization while keeping the sequence short enough for stable pooling.

For each feature channel $c$, masked statistics pooling computes a weighted mean and standard deviation,
\begin{equation}
  \mu_c = \frac{\sum_f w_f z_{c,f}}{\sum_f w_f},
  \qquad
  \sigma_c = \sqrt{\frac{\sum_f w_f(z_{c,f}-\mu_c)^2}{\sum_f w_f} + \varepsilon},
\end{equation}
where $z_{c,f}$ are detector features and $w_f$ are either ground-truth mask weights during training or presence-derived weights during inference. The concatenated vector $[\mu;\sigma]$ is then passed to an MLP for payload recovery.

\begin{lstlisting}[language=Python,caption={Simplified payload decoding with masked statistics pooling.},label={lst:masked_pool}]
def decode_bits(self, features, weights=None):
    z = self.bit_context(features)
    pooled = self.masked_stats_pool(z, weights)
    return self.bit_mlp(pooled)
\end{lstlisting}

This detector is significantly easier to optimize than a per-frame bit decoder and is one of the main reasons the present system converges reliably.

\section{Staged Training Curriculum}\label{sec:curriculum}
\subsection{Why a curriculum was required}
Directly training a localized, low-latency, payload-carrying watermark from scratch led to unstable outcomes in early experiments. Common failure modes were easy to recognize: the detector predicted watermark everywhere, or payload recovery stayed near random even when presence detection looked superficially good. The curriculum breaks this problem into smaller subproblems and avoids changing too many degrees of freedom at once.

\subsection{Implemented stages}
Table~\ref{tab:curriculum} summarizes the stages that were actually used.

\begin{table}[H]
\centering
\caption{Implemented training curriculum used in the current system.}
\label{tab:curriculum}
\begin{tabular}{p{2.0cm}p{4.2cm}p{7.6cm}}
\toprule
\textbf{Stage} & \textbf{Purpose} & \textbf{Main settings} \\
\midrule
Debug stage 0 & verify that the bit channel can learn at all & one-file or tiny-manifest sanity checks, fixed payloads, no attacks, no quality penalty \\
Stage 1 bridge (4-bit) & learn a reliable payload on real speech & full-segment watermark, 4-bit payload, 2~s segments, no negatives, no quality penalty \\
Stage 2A detector calibration & teach the detector to separate watermark from clean audio & full-segment watermark, \texttt{p\_watermark=0.5}, positives and negatives mixed, 4-bit payload \\
Stage 3 localization & introduce localized watermark regions & 4~s segments, dropout embedding with short holes, balanced presence loss, 4-bit payload \\
Stage 4 latency & improve short-prefix payload recovery & stage 3 setup + latency auxiliary loss on prefixes of 0.25/0.5/1.0/2.0~s \\
\bottomrule
\end{tabular}
\end{table}

Across stage transitions, the training script reloads \textbf{model weights only} from the best checkpoint of the previous stage. Optimizer state is intentionally not resumed between stages, because each stage changes the effective objective and the data distribution.

\subsection{Stable stage-3 and stage-4 settings}
The two mature checkpoints discussed in this thesis are:
\begin{itemize}[leftmargin=1.25em]
  \item \textbf{Stage 3 (localized 4-bit)}: dropout embedding with \texttt{ensure\_prefix\_seconds=1.25}, one optional hole of roughly 0.10--0.25~s, and no active quality penalty.
  \item \textbf{Stage 4 (latency 4-bit)}: a fine-tuning stage with \texttt{ensure\_prefix\_seconds=2.0}, very small holes, and an auxiliary prefix-decoding loss applied with probability 0.8.
\end{itemize}
These two stages are not redundant. Stage~3 establishes a stable localized watermark; stage~4 deliberately shifts part of the optimization budget toward early payload recovery.

\subsection{Training objective in the implemented system}
The training loss can be summarized as
\begin{equation}
  \mathcal{L}_{\mathrm{total}} = \lambda_{\mathrm{pres}}\mathcal{L}_{\mathrm{pres}} +
  \lambda_{\mathrm{bit}}\mathcal{L}_{\mathrm{bit}} +
  \lambda_{\mathrm{lat}}\mathcal{L}_{\mathrm{lat}} +
  \lambda_{\mathrm{stft}}\mathcal{L}_{\mathrm{stft}} +
  \lambda_{\mathrm{l2}}\mathcal{L}_{\mathrm{l2}}.
\end{equation}
For the strongest currently reported stage~4 checkpoint, the most important active terms are the presence loss, the bit loss, and the latency loss. The STFT and $\ell_2$ audio-quality terms exist in the code base but were set to zero in the strongest latency-oriented run. This is precisely why perceptual evaluation remains important: the best functional watermark is not automatically the best perceptual watermark.

The latency term reuses detection and payload losses on truncated prefixes:
\begin{lstlisting}[language=Python,caption={Latency auxiliary loss used during stage 4 fine-tuning.},label={lst:latency_loss}]
if latency_enable and random() < latency_p:
    Tchunk = seconds_to_samples(sampled_chunk_seconds)
    y_chunk = y_att[:, :, :Tchunk]
    m_chunk = mask_use[:, :, :Tchunk]
    pres_l, bit_l = compute_detection_losses_only(
        y_chunk, m_chunk, bits, model.detector, cfg
    )
    lat_loss = lambda_pres * pres_l + lambda_bit * bit_l
    loss = loss + latency_weight * lat_loss
\end{lstlisting}

\subsection{Optimizer and validation settings}
The implementation uses AdamW with gradient clipping and fixed random seeds. Mixed precision was tested, but the mature stage~3 and stage~4 runs reported here used \texttt{amp=false}. Validation is performed in two modes after each epoch:
\begin{itemize}[leftmargin=1.25em]
  \item \textbf{clean validation}: watermark applied without extra train-time distortion,
  \item \textbf{attacked validation}: watermark passed through the selected train-time attack operator.
\end{itemize}
The best checkpoint is selected using clean validation bit accuracy.

\section{Evaluation Protocol}
\subsection{Implemented evaluation-time attack families}
Robustness is measured with a dedicated evaluation script that applies fixed black-box transformations outside the differentiable training graph. The retained stage~3 and stage~4 CSVs cover the following attack families:
\begin{itemize}[leftmargin=1.25em]
  \item additive noise,
  \item speed/resampling perturbation,
  \item low-pass filtering,
  \item crop-and-pad,
  \item bit-depth reduction,
  \item MP3 round-trip,
  \item AAC round-trip.
\end{itemize}
For each transformation, the evaluation includes both positive and negative trials. Positive trials attack a watermarked segment; negative trials apply the same channel to the clean segment. This symmetry is important because it lets the system report \textbf{TPR and FPR under the same channel} rather than only on clean negatives.

\subsection{Core metrics}
The implemented benchmark reports:
\begin{itemize}[leftmargin=1.25em]
  \item \textbf{TPR/FPR}: watermark-detection quality,
  \item \textbf{bit accuracy}: fraction of payload bits recovered correctly,
  \item \textbf{message accuracy}: fraction of samples whose full 4-bit payload is recovered,
  \item \textbf{prefix metrics}: TPR and message accuracy at 0.25~s, 0.5~s, 1.0~s, and 2.0~s.
\end{itemize}
The prefix metrics are central to the thesis claim about low-latency verification. They should nevertheless be interpreted carefully: they quantify \emph{watermark-prefix latency}, not the latency of a full assistant stack that would also include streaming ASR, policy logic, networking, and possibly a user-side detector.

\subsection{Objective quality protocol added for RQ1}
Reviewer comment~1 asked for at least a small objective quality assessment comparing clean and watermarked speech. The revised evaluation protocol therefore adds four paired clean-versus-watermarked metrics on top of the robustness benchmark:
\begin{itemize}[leftmargin=1.25em]
  \item \textbf{waveform SNR},
  \item \textbf{segmental SNR},
  \item \textbf{PESQ} for wideband speech quality~[71],
  \item \textbf{STOI} for intelligibility preservation~[72].
\end{itemize}
For a clean signal $x$ and a watermarked signal $y$, the waveform SNR is
\begin{equation}
  \mathrm{SNR}(x,y) = 10 \log_{10}\frac{\|x\|_2^2}{\|x-y\|_2^2}.
\end{equation}
Segmental SNR refines this by averaging over short analysis frames, which makes the metric less dominated by a few high-energy regions. PESQ and STOI complement SNR because a watermark can have acceptable energy-based distortion while still harming perceived speech quality or intelligibility.

To make this protocol reproducible, the accompanying code archive now includes \texttt{quality\_eval.py}, which renders watermarked audio from a checkpoint and writes per-item and aggregate quality statistics. The relevant summary is given in Table~\ref{tab:rq1_quality_metrics}. Importantly, the archive preserved for the present revision contains the stage~3/stage~4 result CSVs but not the stage~4 checkpoint artifact or rendered clean/watermarked audio pairs. For that reason, this revision adds the \emph{protocol and implementation} without inventing numeric PESQ/STOI/SNR values that cannot be recomputed honestly from the retained files.

\begin{table}[H]
\centering
\caption{Objective perceptual metrics added to the revised RQ1 evaluation protocol.}
\label{tab:rq1_quality_metrics}
\begin{tabular}{p{2.6cm}p{7.1cm}p{2.1cm}}
\toprule
\textbf{Metric} & \textbf{What it measures in this thesis} & \textbf{Desired direction} \\
\midrule
SNR & overall energy of watermark residual relative to the clean waveform & higher \\
Segmental SNR & local frame-wise distortion rather than only global energy & higher \\
PESQ & perceptual speech quality under wideband conditions & higher \\
STOI & preservation of intelligibility between clean and watermarked speech & higher \\
\bottomrule
\end{tabular}
\end{table}

The only preserved numeric proxy for stage~4 perturbation in the current archive comes from the demo-rendering path, which reported mean absolute residual magnitudes of roughly 0.16--0.18 and maximum residuals hitting the configured 0.18 ceiling on randomly sampled examples. This is \emph{not} a substitute for PESQ/STOI, but it does support an important qualitative conclusion: the strongest latency-focused checkpoint should be interpreted as a strong \emph{functional} watermark rather than as a watermark already optimized for strict inaudibility.

\subsection{Additional channel diversity added for RQ2}
Reviewer comment~2 requested a broader channel benchmark than the original offline attack suite. A full VoIP or room-acoustic study is still outside the current project scope, but the revised evaluation code now adds two practical channel operators:
\begin{itemize}[leftmargin=1.25em]
  \item \textbf{low-bitrate Opus round-trip} (e.g., 12 and 24~kb/s), which approximates modern conversational coding better than MP3/AAC alone~[74],
  \item a \textbf{PSTN-style phone-band proxy}, which applies an 8~kHz narrowband path, a 300--3400~Hz passband, and optional $\mu$-law companding in the spirit of legacy telephony~[75].
\end{itemize}
This extension is intentionally modest, but it is methodologically important. It moves the benchmark one step closer to channels encountered in practical assistants and call-like settings. As with the quality protocol above, the retained archive did not include the stage~4 checkpoint needed to execute these newly added rows retroactively. Nevertheless, the code path is now present in \texttt{robustness\_eval.py}, so the additional benchmark can be produced directly once the checkpoint artifact is restored.

\section{Empirical Results and Discussion}
\subsection{Stage 3: stable localized baseline}
The clean \texttt{stage3\_localize.csv} row shows that stage~3 already establishes a strong baseline on LibriSpeech. On clean evaluation, the localized stage achieves:
\begin{itemize}[leftmargin=1.25em]
  \item TPR = 1.000,
  \item FPR = 0.010,
  \item bit accuracy = 0.9863,
  \item full message accuracy = 0.975,
  \item message accuracy at 0.25/0.5/1.0/2.0~s = 0.015 / 0.105 / 0.265 / 0.440.
\end{itemize}
These numbers are useful for two reasons. First, they show that localization-compatible training does not destroy payload recovery. Second, they reveal the central limitation that motivated stage~4: stage~3 is already a strong \emph{full-segment} watermark, but it is still weak on very short prefixes.

\subsection{Stage 4: latency-aware fine-tuning}
The clean \texttt{stage4\_latency.csv} row confirms that stage~4 shifts the operating point in the intended direction. Table~\ref{tab:stage3_stage4_clean_exact} compares the clean stage~3 and stage~4 checkpoints using the values preserved in the uploaded CSV files.

\begin{table}[H]
\centering
\caption{Clean stage~3 versus stage~4 results from the retained CSV artifacts.}
\label{tab:stage3_stage4_clean_exact}
\begin{tabular}{lccc}
\toprule
\textbf{Metric} & \textbf{Stage 3} & \textbf{Stage 4} & \textbf{Change} \\
\midrule
TPR & 1.000 & 1.000 & 0.000 \\
FPR & 0.010 & 0.010 & 0.000 \\
Bit accuracy & 0.9863 & 0.9825 & $-0.0038$ \\
Full message accuracy & 0.975 & 0.960 & $-0.015$ \\
Message accuracy @ 0.25~s & 0.015 & 0.635 & $+0.620$ \\
Message accuracy @ 0.5~s  & 0.105 & 0.675 & $+0.570$ \\
Message accuracy @ 1.0~s  & 0.265 & 0.690 & $+0.425$ \\
Message accuracy @ 2.0~s  & 0.440 & 0.785 & $+0.345$ \\
\bottomrule
\end{tabular}
\end{table}

\begin{figure}[H]
\centering
\includegraphics[width=0.92\textwidth]{stage3_stage4_prefix_msgacc.png}
\caption{Clean prefix message accuracy for stage~3 and stage~4. Stage~4 deliberately trades a small amount of full-length message accuracy for a much earlier payload recovery curve.}
\label{fig:stage3_stage4_prefix_msgacc}
\end{figure}

The practical interpretation is clear. Stage~4 does \emph{not} produce a uniformly better checkpoint in every metric; instead, it makes a targeted trade-off that is highly relevant to interactive use. The false-positive rate stays unchanged at 0.01 and full clean message accuracy decreases only slightly, but short-prefix decoding improves dramatically. This is exactly the behavior that the latency auxiliary loss was designed to encourage.

The same pattern is visible beyond the clean row. Averaged across all attacks shared by the retained stage~3 and stage~4 CSVs, stage~4 improves message accuracy by about +0.508 at 0.25~s, +0.545 at 0.5~s, +0.349 at 1.0~s, and +0.294 at 2.0~s, while sacrificing only about 0.024 in full-length message accuracy. In other words, stage~4 meaningfully improves early verification not only in the easiest case but across the entire retained robustness sweep.

\subsection{Robustness under the implemented benchmark}
Table~\ref{tab:stage4_selected_attacks} summarizes several representative stage~4 attack conditions drawn directly from \texttt{results/stage4\_latency.csv}.

\begin{table}[H]
\centering
\scriptsize
\setlength{\tabcolsep}{4pt}
\caption{Selected stage~4 robustness results from the retained benchmark.}
\label{tab:stage4_selected_attacks}
\begin{tabular}{lcccccc}
\toprule
\textbf{Attack} & \textbf{TPR} & \textbf{FPR} & \textbf{Bit acc.} & \textbf{Msg acc.} & \textbf{Msg acc. @0.5~s} & \textbf{Msg acc. @2.0~s} \\
\midrule
Clean & 1.000 & 0.010 & 0.9825 & 0.9600 & 0.675 & 0.785 \\
Noise (SNR 5 dB) & 1.000 & 0.010 & 0.9800 & 0.9550 & 0.690 & 0.760 \\
Speed 0.9$\times$ & 1.000 & 0.015 & 0.9450 & 0.8050 & 0.620 & 0.705 \\
Low-pass 3 kHz & 1.000 & 0.010 & 0.9825 & 0.9600 & 0.665 & 0.760 \\
Crop 1.0 s & 0.995 & 0.005 & 0.9196 & 0.7136 & 0.665 & 0.765 \\
MP3 32 kb/s & 1.000 & 0.010 & 0.9850 & 0.9750 & 0.675 & 0.790 \\
AAC 48 kb/s & 1.000 & 0.010 & 0.8513 & 0.4600 & 0.640 & 0.660 \\
\bottomrule
\end{tabular}
\end{table}

Three conclusions matter most for \textbf{RQ2}. First, the watermark is robust to several common offline distortions: additive noise, low-pass filtering, and MP3 produce little or no degradation in TPR, and full message accuracy remains close to the clean case. Second, some timing and editing perturbations are distinctly harder: speed 0.9$\times$ and crop-and-pad both reduce full message accuracy, even if the detector usually still identifies watermark presence. Third, transform coding is not uniformly benign. MP3 is surprisingly easy for the current model, but AAC at 48~kb/s is a clear weak point, dropping full message accuracy from 0.96 to 0.46. This is a large enough gap to matter for deployment-oriented claims.

These results also help interpret the requested channel-diversity extension. The existing \texttt{lowpass\_3000} row shows that pure bandwidth truncation alone is not catastrophic in the current checkpoint; by contrast, the AAC row shows that coding artifacts can be much more destructive than a simple spectral bottleneck. That observation is precisely why the added Opus and PSTN-style operators are useful: they will determine whether practical conversational channels behave more like the relatively mild low-pass case or the harsher codec case.

\subsection{Candidate middleware interaction and practical prefix verification}
Reviewer comment~3 asked for a clearer explanation of how a future user-side detector would coexist with the watermark branch and how prefix-aware watermark verification interacts with streaming ASR latency. Figure~\ref{fig:middleware_watermark_detector_diagram} summarizes the intended interaction.

\begin{figure}[H]
\centering
\includegraphics[width=\textwidth]{middleware_watermark_detector_diagram.png}
\caption{Conceptual middleware interaction between the implemented assistant-output watermark branch and the planned user-side detector branch. The two branches are logically separate but feed a shared trust state consumed by the dialogue policy.}
\label{fig:middleware_watermark_detector_diagram}
\end{figure}

At the systems level, the two branches solve complementary problems. The \textbf{assistant-to-user branch} verifies provenance of assistant audio: the watermark encoder produces a known 4-bit payload, and a verifier running on the received audio stream estimates watermark presence and payload consistency on successive prefixes. The \textbf{user-to-assistant branch} would inspect incoming microphone or uplink audio for replay or deepfake evidence before high-risk actions are accepted. The two branches do not need to share a feature extractor or a joint model. What they \emph{do} share is a \textbf{middleware trust state}: assistant-audio provenance and user-audio integrity become two inputs to the same dialogue policy layer.

The interaction with ASR latency is subtle but important. Watermark decoding does not require the ASR transcript; it can run directly on the audio prefix while playback is still ongoing. ASR, in turn, produces partial hypotheses that may arrive before watermark confidence is fully mature. In practice this means the pipeline should not be interpreted as \emph{serial} (``finish watermark verification, then start ASR''). A better design is \emph{parallel}: ASR produces hypotheses continuously, the watermark verifier updates provenance confidence at 0.25/0.5/1.0/2.0~s checkpoints, and the dialogue manager decides whether the current trust state is already sufficient for the action under consideration.

Table~\ref{tab:prefix_policy} shows one conservative operational policy derived from the clean stage~4 prefix metrics. The model-level detection thresholds remain $\tau=0.5$ and $\rho=0.05$, but the \emph{middleware} can apply an additional risk-aware rule: early windows may set a provisional provenance flag, whereas irreversible or security-sensitive actions should wait for either a longer prefix or repeated agreement across adjacent windows.

\begin{table}[H]
\centering
\caption{Practical interpretation of stage~4 clean prefix metrics for middleware gating.}
\label{tab:prefix_policy}
\begin{tabular}{p{1.8cm}p{1.6cm}p{2.0cm}p{8.0cm}}
\toprule
\textbf{Prefix} & \textbf{TPR} & \textbf{Msg acc.} & \textbf{Suggested middleware interpretation} \\
\midrule
0.25~s & 1.000 & 0.635 & Strong early presence cue, but payload evidence is still too immature for a final trust decision in a high-risk setting. Suitable for a provisional \emph{likely authentic} flag only. \\
0.5~s & 0.995 & 0.675 & Useful point for tentative provenance. Low-risk UI cues could already surface, but security-relevant commitments should ideally wait for another consistent window or a longer prefix. \\
1.0~s & 1.000 & 0.690 & Reasonable operating point for low-risk confirmation if repeated windows agree, because presence is stable and payload reliability is materially better than in stage~3. \\
2.0~s & 1.000 & 0.785 & Best confirmation point in the current checkpoint for final provenance gating; still not perfect, but substantially stronger than earlier prefixes and appropriate for conservative policy thresholds. \\
\bottomrule
\end{tabular}
\end{table}

This table clarifies how prefix-aware verification interacts with conversational timing. The watermark branch need not delay every ASR hypothesis; instead, it informs \emph{when the system may trust and act on those hypotheses}. In a practical assistant, low-risk operations could proceed under provisional trust, while sensitive actions (for example, payment authorization or device control) would require a later and more stable provenance state.

\subsection{Direct status of RQ1--RQ4}
The revised results can now be mapped back to the research questions more precisely.
\begin{itemize}[leftmargin=1.25em]
  \item \textbf{RQ1 (Perceptual impact): partially answered.} The revision now defines and implements a proper objective PESQ/STOI/SNR protocol, but the retained archive does not include the stage~4 checkpoint needed to fill in the numeric table retroactively. The preserved residual-magnitude evidence nevertheless indicates that the best latency-focused checkpoint is not yet optimized for imperceptibility.
  \item \textbf{RQ2 (Robustness on the implemented benchmark): answered within scope.} The watermark survives multiple black-box transformations with high TPR and strong message recovery, especially under noise, low-pass filtering, and MP3. AAC and some timing-related perturbations remain the main weaknesses. The benchmark scope has also been widened in code to include Opus and a phone-band proxy for the next rerun.
  \item \textbf{RQ3 (Localized, low-latency detection): answered within scope.} Stage~3 provides the stable localized baseline, and stage~4 clearly improves short-prefix payload recovery without increasing the clean false-positive rate.
  \item \textbf{RQ4 (Integration status): partially answered.} The thesis still does not measure a full end-to-end middleware, but it now explains concretely how the assistant-output verifier and a future user-side detector would interact through a shared trust state and how policy decisions can be aligned with ASR hypothesis timing.
\end{itemize}

\section{Qualitative Audio Inspection and Demo Examples}
To complement aggregate metrics, the code base includes scripts that render random examples from a manifest, save both original and watermarked waveforms, and run inference on both. The resulting bundle contains:
\begin{itemize}[leftmargin=1.25em]
  \item \texttt{original.wav},
  \item \texttt{watermarked.wav},
  \item \texttt{delta.wav},
  \item \texttt{delta\_normalized.wav},
  \item \texttt{delta\_xgain.wav},
  \item \texttt{meta.json},
  \item summary CSV/JSON outputs with detector predictions.
\end{itemize}
The retained stage~4 demo summary reported that, on five randomly sampled files, every clean example was classified as \emph{not watermarked}, while every watermarked example was classified as \emph{watermarked} and decoded with 100\% payload accuracy. This is a useful sample-level confirmation that the detector behaves sensibly on individual files. At the same time, the same demo path also revealed why RQ1 remains only partially answered: the residual energy was visibly non-negligible, and the maximum perturbation repeatedly reached the configured ceiling. In other words, the present checkpoint is strongest as a provenance mechanism, not yet as a perceptually optimized watermark.

\section{Implementation Notes and Reproducibility}
The current code base is organized around a small number of explicit components:
\begin{itemize}[leftmargin=1.25em]
  \item \texttt{watermark/model.py}: encoder, detector, and combined network,
  \item \texttt{watermark/runtime.py}: model construction, embedding helpers, and detector helpers,
  \item \texttt{train.py}: staged training loop,
  \item \texttt{eval.py}: clean/attacked evaluation,
  \item \texttt{robustness\_eval.py}: black-box attack sweep plus prefix probes,
  \item \texttt{quality\_eval.py}: paired clean-versus-watermarked quality metrics,
  \item \texttt{scripts/render\_watermark\_examples.py}: qualitative listening examples.
\end{itemize}
The revised robustness path additionally supports low-bitrate Opus and a PSTN-style phone-band proxy, which makes the repository more directly aligned with the reviewer request for modest channel realism.

A representative stage~4 training command is:
\begin{lstlisting}[language=bash,caption={Example stage~4 training command.},label={lst:stage4_cmd}]
python train.py \
  --config configs/stage4_latency_4bit.yaml \
  --train_manifest manifests/train_2k.csv \
  --val_manifest manifests/val_200.csv \
  --resume checkpoints_stage3_localize_4bit/best_clean_bitacc.pt
\end{lstlisting}

Robustness and latency evaluation are run as follows:
\begin{lstlisting}[language=bash,caption={Example robustness and latency evaluation command.},label={lst:robustness_cmd}]
python robustness_eval.py \
  --config configs/stage4_latency_4bit.yaml \
  --checkpoint checkpoints_stage4_latency_4bit/best_clean_bitacc.pt \
  --manifest manifests/librispeech_test.csv \
  --out results/stage4_latency.csv \
  --latency
\end{lstlisting}

Objective quality evaluation can be reproduced with:
\begin{lstlisting}[language=bash,caption={Example objective quality evaluation command.},label={lst:quality_eval_cmd}]
python quality_eval.py \
  --config configs/stage4_latency_4bit.yaml \
  --checkpoint checkpoints_stage4_latency_4bit/best_clean_bitacc.pt \
  --manifest manifests/librispeech_test.csv \
  --out_csv results/stage4_quality_rows.csv \
  --out_json results/stage4_quality_summary.json
\end{lstlisting}

\section{Current Limitations and Next Step}
The present prototype has four important limitations.
\begin{enumerate}[leftmargin=1.25em]
  \item It is still evaluated only on speech from LibriSpeech rather than on a broader multi-domain audio benchmark.
  \item The strongest checkpoint is optimized primarily for detection, localization, and prefix-aware decoding; it has not yet been convincingly optimized for imperceptibility.
  \item The user-side deepfake/replay detector remains a future extension rather than an implemented experimental branch.
  \item Full end-to-end middleware latency is still not measured directly; only watermark-prefix behavior is quantified.
\end{enumerate}

These limitations lead to a clear next step. First, the preserved stage~4 checkpoint should be rerun through the newly added quality and channel-diversity scripts so that PESQ/STOI/SNR, Opus, and phone-band rows become part of the main experimental table rather than only the protocol. Second, a quality-oriented fine-tuning stage should reduce \texttt{max\_delta} and activate small STFT/$\ell_2$ penalties to study the fidelity-versus-recovery trade-off explicitly. Third, once that checkpoint is stable, the planned user-side detector branch can be integrated into a shared middleware for a true end-to-end latency study. This is therefore an honest and useful stopping point for the current thesis: the prototype already demonstrates a localized, low-FPR, low-latency assistant-output watermark, while the remaining work concerns perceptual refinement, broader channel realism, and bilateral systems integration.

\FloatBarrier

\section*{References}
\small
\noindent [1] [15] (PDF) On the Security and Privacy Challenges of Virtual Assistants. \url{https://www.researchgate.net/publication/350413576_On_the_Security_and_Privacy_Challenges_of_Virtual_Assistants}\par
\noindent [2] [6] [7] [8] [13] [17] [18] [38] [39] [40] [41] [42] [43] [44] [45] [48] [2308.14970] Audio Deepfake Detection: A Survey. \url{https://ar5iv.labs.arxiv.org/html/2308.14970}\par
\noindent [3] [16] usenix.org. \url{https://www.usenix.org/system/files/usenixsecurity24-jiang-nan.pdf}\par
\noindent [4] [5] [12] [23] [24] [29] [30] [31] [49] usenix.org. \url{https://www.usenix.org/system/files/usenixsecurity25-zong.pdf}\par
\noindent [9] [11] [19] [20] [21] [22] Audio Watermarking Won't Solve the Real Dangers of AI Voice Manipulation -- Center for Data Innovation. \url{https://datainnovation.org/2024/10/audio-watermarking-wont-solve-the-real-dangers-of-ai-voice-manipulation/}\par
\noindent [10] [28] [35] [36] [37] [46] [47] [50] [51] [52] [53] [54] FakeMark: Deepfake Speech Attribution With Watermarked Artifacts. \url{https://arxiv.org/html/2510.12042v2}\par
\noindent [14] Proactive Detection of Voice Cloning with Localized Watermarking. \url{https://ai.meta.com/research/publications/proactive-detection-of-voice-cloning-with-localized-watermarking/}\par
\noindent [25] [26] [27] [2401.17264] Proactive Detection of Voice Cloning with Localized Watermarking. \url{https://arxiv.org/abs/2401.17264}\par
\noindent [32] [33] [34] [2504.14832] Protecting Your Voice: Temporal-aware Robust Watermarking. \url{https://arxiv.org/abs/2504.14832}\par
\noindent [55] SynthID --- Google DeepMind. \url{https://deepmind.google/models/synthid/}\par
\noindent [56] Proactive Detection of Voice Cloning with Localized Watermarking (AudioSeal). \url{https://arxiv.org/abs/2401.17264}\par
\noindent [57] AudioMarkBench: Benchmarking Robustness of Audio Watermarking. \url{https://arxiv.org/abs/2406.06979}\par
\noindent [58] LibriSpeech ASR corpus (OpenSLR SLR12). \url{https://www.openslr.org/12/}\par
\noindent [70] ITU-T Rec. P.800.1: Mean Opinion Score (MOS) terminology. \url{https://www.itu.int/rec/dologin_pub.asp?id=T-REC-P.800.1-200303-S!!PDF-E&lang=e&type=items}\par
\noindent [71] ITU-T Rec. P.862: Perceptual evaluation of speech quality (PESQ). \url{https://www.itu.int/rec/T-REC-P.862/}\par
\noindent [72] Taal et al.\ (2011): Short-Time Objective Intelligibility (STOI). \url{https://sps.ewi.tudelft.nl/pubs/Taal2011_1.pdf}\par
\noindent [73] AudioSeal GitHub repository. \url{https://github.com/facebookresearch/audioseal}\par
\noindent [74] RFC 6716: Definition of the Opus Audio Codec. \url{https://datatracker.ietf.org/doc/html/rfc6716}\par
\noindent [75] ITU-T Rec. G.711: Pulse code modulation (PCM) of voice frequencies. \url{https://www.itu.int/rec/T-REC-G.711}\par
